\exercisesection{正则局部环}

\begin{enumerate}
\def\labelenumi{\arabic{enumi})}
\item
  设 A 为一个环，S 为 A 的一个乘法子集。证明有 \(\mathrm{dh}(\mathrm{S}^{-1}\mathrm{A}) \leqslant \mathrm{dh}(\mathrm{A})\)（A, X, p.138, 定义 2）。（取 \(S^{-1}A\)-模的投射 A-模解并将其局部化。）
\item
  设 A 为诺特环，B 为忠实平坦的 A-代数。则有 \(\mathrm{dh(A)}\leqslant \mathrm{dh(B)}\)。
\end{enumerate}

3) 设 A 为诺特局部环。我们拟证明，A 为正则环当且仅当 dh(A) 有限。

\begin{enumerate}
\def\labelenumi{\alph{enumi})}
\tightlist
\item
  假设 dh(A) 有限。为证明 A 为正则环，使用 A, X, p.208, 习题 13。
\item
  假设 A 为正则环。设 x 是生成 m\_A 的完全相交族。则与 x 相关的 Koszul 复形是由自由 A-模构成的 A/m\_A 的有限解。然后应用 A, X, p.203, 习题 16。
\end{enumerate}

\begin{enumerate}
\def\labelenumi{\arabic{enumi})}
\setcounter{enumi}{3}
\item
  设 A 为正则诺特局部环。则有 \(\dim(A) = \mathrm{dh}(A)\)（为证明 \(\mathrm{dh}(A) \leqslant \dim(A)\)，使用与参数系相关的 Koszul 复形。然后确定群 \(\operatorname{Ext}_{A}^{i}(\kappa_{A}, \kappa_{A})\)。）
\item
  设 A 为正则诺特局部环，p 为 A 的一个素理想。则 \(A_{p}\) 为正则环（使用习题 1 和习题 3）。
\item
  若对于 A 的每个素理想 p，局部环 \(A_{p}\) 都是正则环，则称诺特环 A 为正则环。
\end{enumerate}

\begin{enumerate}
\def\labelenumi{\alph{enumi})}
\item
  设 A 为诺特环。则 A 为正则环当且仅当对于 A 的每个极大理想 m，\(A_{m}\) 都是正则环（使用习题 5）。
\item
  设 A 为诺特环。证明 A 为有限维正则环当且仅当 dh(A) 有限（使用习题 4）。
\item
  证明 p.83 习题 13 中描述的环是无限维正则环。
\item
  设 A 为正则诺特环，S 为一个乘法子集。证明 \(\mathbf{S}^{-1}\mathbf{A}\) 为正则环。
\item
  设 A 为正则诺特环。证明 A{[}X{]} 为正则环。（通过局部化归约到 dh(A) 有限的情形。使用 A, X, p. 143, 定理 1 及 b)。）
\end{enumerate}

\(^{1}\) 关于此习题的更多细节，可参阅 R. Stanley，《分次代数的 Hilbert 函数》，Advances in Math., 28 (1978), 第 57–83 页。

\begin{enumerate}
\def\labelenumi{\arabic{enumi})}
\setcounter{enumi}{6}
\tightlist
\item
  设 I 为含有可数共尾子集的定向集，A 为一个环（不必交换）。我们拟证明，以 I 为指标的投射（左）A-模的每个归纳极限的投射维数都 \(\leqslant 1\)（A, X, p.~134）。
\end{enumerate}

\begin{enumerate}
\def\labelenumi{\alph{enumi})}
\tightlist
\item
  设 \(((\mathbf{Q}_i)_{i\in \mathbf{I}},(\varphi_{i,j})_{i < j})\) 为投射 A-模的归纳系统，且 \(Q = \varinjlim_{I} Q_i\)。为证明 \(Q\) 的投射维数 \(\leqslant 1\)，归约到 \(I = \mathbf{N}\) 的情形。此时有正合序列
\end{enumerate}

\[
0 \longrightarrow \bigoplus_ {n \in \mathbf {N}} Q _ {n} \stackrel {u} {\longrightarrow} \bigoplus_ {n \in \mathbf {N}} Q _ {n} \longrightarrow Q \longrightarrow 0
\]

其中对于每个 \(n \in \mathbf{N}\) 及每个 \(x \in \mathbf{Q}_n\)，有

\[
u (x) = \varphi_ {n, n + 1} (x) - x.
\]

\begin{enumerate}
\def\labelenumi{\alph{enumi})}
\setcounter{enumi}{1}
\tightlist
\item
  设 \(((\mathbf{P}_i)_{i\in I}, (\varphi_{i,j})_{i < j})\) 为投射 A-模的归纳系统，使得对于每个对 \((i,j)\)（其中 \(i < j\)），\(\varphi_{i,j}\) 都是单射，且模 \(\operatorname{Coker}(\varphi_{i,j})\) 是投射的。证明模 \(\mathbf{P} = \varinjlim_{I} \mathbf{P}_i\) 是投射的。（设 \(0 \to \mathbf{M}' \to \mathbf{M} \to \mathbf{M}'' \to 0\) 是 A-模的正合序列。则
\end{enumerate}

\[
0 \to \operatorname{Hom} _ {\mathbf {A}} (\mathrm{P} _ {i}, \mathrm{M} ^ {\prime}) \to \operatorname{Hom} _ {\mathbf {A}} (\mathrm{P} _ {i}, \mathrm{M}) \to \operatorname{Hom} _ {\mathbf {A}} (\mathrm{P} _ {i}, \mathrm{M} ^ {\prime \prime}) \to 0
\]

这是交换群的射影系统的正合序列，且射影系统 \(i \mapsto \mathrm{Hom}_{\mathbf{A}}(\mathbf{P}_{i}, \mathbf{M}')\) 对离散拓扑具有 Mittag-Leffler 性质（TG, II, p.18）。由此可推出射影极限序列是正合的。）

\begin{enumerate}
\def\labelenumi{\alph{enumi})}
\setcounter{enumi}{2}
\tightlist
\item
  设 \((\mathbf{Q}_i)_{i\in I}\) 为投射 A-模的归纳系统。对于每个 \(i\in \mathbf{I}\)，令 \(\mathbf{R}_i = \bigoplus_{j\leqslant i}\mathbf{Q}_j\)，并且对于 \(i < i'\)，令 \(\psi_{i,i'}: \mathbf{R}_i \to \mathbf{R}_{i'}\) 为典范单射。定义归纳系统的正合序列 \(0 \to (\mathbf{P}_i) \to (\mathbf{R}_i) \to (\mathbf{Q}_i) \to 0\)，使得 \((\mathbf{P}_i)\) 满足 \(b\) 的假设。由此推出 \(a\) 的另一种证明。）
\end{enumerate}

\begin{enumerate}
\def\labelenumi{\arabic{enumi})}
\setcounter{enumi}{7}
\tightlist
\item
  设 A 为一个环，n 为满足 n ≥ 1 的整数。证明下列条件等价：
\end{enumerate}

\begin{enumerate}
\def\labelenumi{(\roman{enumi})}
\item
  有 dh(A) ≤ n。
\item
  对每个循环 A-模 M，都有 \(\mathrm{dp}_{\mathbf{A}}(\mathbf{M}) \leqslant n\)。
\item
  对 A 的每个理想 I，都有 \(\mathrm{dp}_{\mathbf{A}}(\mathbf{I}) \leqslant n - 1\)。
\end{enumerate}

(使用 A, X, p.138, 命题 4 证明等价性 (i) \(\Leftrightarrow\) (ii)，并使用 A, X, p.93, 命题 11 证明 (i) \(\Leftrightarrow\) (iii)。)

\begin{enumerate}
\def\labelenumi{\arabic{enumi})}
\setcounter{enumi}{8}
\tightlist
\item
  设 \(\Gamma\) 为全序交换群，\(\Gamma^{+}\) 为 \(\Gamma\) 的正元素集。我们以 (GOD) 表示下列性质：
\end{enumerate}

(GOD)：每个主子集 M ⊂ Γ⁺（VI, § 3, no. 5, 定义 2）都有一个关于逆序共尾的可数子集。

\begin{enumerate}
\def\labelenumi{\alph{enumi})}
\item
  证明若 \(\Gamma\) 的高度为 1，则满足性质 (GOD)。
\item
  对每个整数 n，证明存在一个具有性质 (GOD) 且高度为 n 的全序群 \(\Gamma\)（取赋予字典序的 \(\Gamma = \mathbf{Z}^{n}\)）。
\end{enumerate}

10) 设 A 为赋值环，其值的有序群 \(\Gamma\) 具有性质 (GOD)。a) 证明 A 的同调维数至多为 2。（注意 A 的每个理想都是主理想的可数归纳极限，并使用习题 7 和 8 的结果。）

\begin{enumerate}
\def\labelenumi{\alph{enumi})}
\setcounter{enumi}{1}
\item
  每个 Krull 维数为 1 的赋值环的同调维数都 \(\leqslant 2\)。它的同调维数为 2 当且仅当它不是诺特环。（使用习题 9, a) 以及 A, X, p.208, 习题 12。）
\item
  对每个整数 \(n \geqslant 1\)，存在 Krull 维数等于 n 且同调维数等于 2 的环。（使用习题 9, b) 以及 A, X, p.208, 习题 12。）
\end{enumerate}

\begin{enumerate}
\def\labelenumi{\arabic{enumi})}
\setcounter{enumi}{10}
\item
  每个正则诺特局部环都是唯一分解环。（使用 VII, § 4, n° 7, 推论 3。）
\item
  设 A 为一个环，a 为 A 的有限生成理想。
\end{enumerate}

\begin{enumerate}
\def\labelenumi{\alph{enumi})}
\item
  证明若 A 存在理想 b 使得 a = a·b，则存在 \(b \in b\) 使得 \((1 - b)a = (0)\)。
\item
  证明若 \(a^{2} = a\)，则存在元素 \(a \in a\) 使得 \(a^{2} = a\) 且 a = Aa。
\item
  假设 a 是极大理想，且 A-模 \(a/a^{2}\) 可由 n 个元素生成。证明理想 a 可由 \(n + 1\) 个元素生成。（取 a 中的元素 \(x_{1}, \ldots, x_{n}\)，它们在 \(a/a^{2}\) 中的像生成后者。记 x 为由 \(\{x_{1}, \ldots, x_{n}\}\) 生成的 A 的理想；注意在 A/x 中有 \((\mathfrak{a}/\mathfrak{x})^{2} = (\mathfrak{a}/\mathfrak{x})\)。）
\end{enumerate}

\begin{enumerate}
\def\labelenumi{\arabic{enumi})}
\setcounter{enumi}{12}
\tightlist
\item
  设 \(p\) 为素数，\(f\) 为满足 \(> 0\) 的整数。置 \(L = \mathbf{Q}(\zeta)\)，其中 \(\zeta\) 是 1 的一个本原 \(p^f\) 次根，并令 \(B = \mathbf{Z}[\zeta]\) 为 L 中由 \(\zeta\) 生成的子环。证明 B 是 Z 在 L 中的整闭包。（通过判别式计算可证明，\(B\left[\frac{1}{p}\right]\) 是 \(\mathbf{Z}\left[\frac{1}{p}\right]\) 的整闭包（V, § 1, no. 6, 引理 3）；记 \(\overline{\mathbf{B}}\) 为 \(\mathbf{B}\) 的整闭包，
\end{enumerate}

由此可推出 \(\overline{\mathbf{B}}\subset\bigcap_{q}\mathbf{Z}_{q}[\zeta]\)，其中 \(q\) 遍历素数，从而 \(\overline{\mathbf{B}}\subset\mathbf{Z}[\zeta].\)

\begin{enumerate}
\def\labelenumi{\arabic{enumi})}
\setcounter{enumi}{13}
\tightlist
\item
  设 \(\rho: A \to B\) 为诺特环的同态，使 \(B\) 成为忠实平坦的 \(A\)-模。
\end{enumerate}

\begin{enumerate}
\def\labelenumi{\alph{enumi})}
\item
  若 B 为正则环，则 A 为正则环（参见 p. 94, 习题 2 和 3）。
\item
  若 A 为正则环，且对于 A 的每个极大理想 m，B/mB 都是正则环，则 B 为正则环。（通过局部化归约到 A、B 为局部环且 ρ 为局部同态的情形。此时有 dim(B) = dim(A) + dim(B/mB)（§ 3, no. 4, 命题 7 的推论 1）。为 B 构造一个含有 dim(B) 项的正则序列。）
\end{enumerate}

\begin{enumerate}
\def\labelenumi{\arabic{enumi})}
\setcounter{enumi}{14}
\item
  设 X 为诺特空间，U 为 X 的子集。若对于 X 中每个与 U 相交的不可约闭子集 Y，U ∩ Y 都含有 Y 的一个非空开子集，则 U 在 X 中是开的。（反设 U 不是开集，考虑满足 Z ∩ U 不是开集的 X 的极小闭集 Z ⊂ X。）
\item
  设 A 为诺特环。Spec(A) 的奇异轨迹记为 Sing(A)，它是满足 \(p \in \text{Spec}(A)\) 且 \(A_p\) 不是正则环的 p 的集合。置
\end{enumerate}

\[
\operatorname{Reg} (\mathrm{A}) = \operatorname{Spec} (\mathrm{A}) - \operatorname{Sing} (\mathrm{A}).
\]

证明下列条件等价：

\begin{enumerate}
\def\labelenumi{(\roman{enumi})}
\item
  对 A 的每个商环 B，Reg(B) 在 Spec(B) 中是开集。
\item
  对 A 的每个整环商环 B，Reg(B) 都是 Spec(B) 的非空开子集。
\item
  对 A 的每个整环商环 B，Reg(B) 都含有 Spec(B) 的一个非空开子集。
\end{enumerate}

（为证明 (i) \(\Rightarrow\) (ii)，注意 \((0)\in\mathrm{Reg}(\mathbf{B})\)。为证明 (iii) \(\Rightarrow\) (i)，证明对于每个 \(\mathfrak{p}\in\mathrm{Spec}(\mathbf{A})\)，使得 \(\mathrm{V}(\mathfrak{p})\cap\mathrm{Reg}(\mathbf{A})\neq\emptyset\)，环 \(\mathbf{A}_{\mathfrak{p}}\) 是正则环（p.94, 习题 5）。使用 (iii) 并构造完全相交序列，证明 \(\mathrm{V}(\mathfrak{p})\cap\mathrm{Reg}(\mathbf{A})\) 含有 \(\mathrm{V}(\mathfrak{p})\) 的一个非空开子集。使用习题 15 得出结论。）

\begin{enumerate}
\def\labelenumi{\arabic{enumi})}
\setcounter{enumi}{16}
\tightlist
\item
  设 \(\rho: A \to B\) 为整诺特环的单同态，使 \(B\) 成为有限型 \(A\)-模。
\end{enumerate}

\begin{enumerate}
\def\labelenumi{\alph{enumi})}
\item
  假设 \(\operatorname{Reg}(\mathbf{B})\) 含有 \(\operatorname{Spec}(\mathbf{B})\) 的一个非空开子集。证明 \(\operatorname{Reg}(\mathbf{A})\) 含有 \(\operatorname{Spec}(\mathbf{A})\) 的一个非空开子集（使用 II, § 5, no. 1, 命题 2 的推论，归约到 \(\mathbf{B}\) 在 \(\mathbf{A}\) 上平坦的情形；然后在 \(\mathbf{B}\) 为正则环的情形，对适当的元素 \(x \in \mathbf{A}\) 局部化。使用习题 14, a) 得出结论。）
\item
  假设 \(\operatorname{Reg}(\mathbf{A})\) 含有 \(\operatorname{Spec}(\mathbf{A})\) 的一个非空开子集，且 \(\mathbf{K}'\)（\(\mathbf{B}\) 的分式域）是 \(\mathbf{K}\)（\(\mathbf{A}\) 的分式域）的可分扩张。证明 \(\operatorname{Reg}(\mathbf{B})\) 含有 \(\operatorname{Spec}(\mathbf{B})\) 的一个非空开子集。（取 \(\mathbf{Z} \in \mathbf{K}'\) 使得 \(\mathbf{K}' = \mathbf{K}(\mathbf{Z})\)（A, V, p.39），令 \(\mathrm{P}(\mathrm{X}) \in \mathrm{K}[\mathrm{X}]\) 为 Z 的最小多项式。用 A 中适当的元素 f 将 A 替换为环 \(A_{f}\)，证明可以归约到证明 Spec(B) 在下列情形中具有所要求的性质：
\end{enumerate}

\(\alpha)\) \(\mathbf{Z}\in \mathbf{B}\) 且 A 为正则环；

\(\beta)\) 有 \(\mathrm{P}(\mathrm{X})\in \mathrm{A}[\mathrm{X}]\)；

\(\gamma)\) 有 \(\mathbf{B} = \mathbf{A}[\mathbf{X}] / (\mathbf{P}(\mathbf{X}))\)（II, § 5, no. 1, 命题 2）；

δ) P'(Z) 在 B 中可逆（注意 P'(Z) ≠ 0 在 K' 中，因为 K' 在 K 上可分。然后使用 II, § 5, no. 1, 命题 3 的推论）。

  证明在此情形下，对于 A 的每个极大理想 m，B/mB 是一个平展 (A/m)-代数（A, V, p. 32），因而是正则环。使用习题 14, b) 得出结论。

18) 设 \(k\) 为一个域，A 为有限型 \(k\)-代数。证明 \(\operatorname{Reg}(A)\) 是 \(\operatorname{Spec}(A)\) 的开子集。（使用习题 16，归约到 A 为整环时证明 \(\operatorname{Reg}(A)\) 含有一个非空开子集。使用正规化引理（V, § 3, no. 1, 定理 1）和习题 17, a)，归约到 A 是 \(k[X_1, ..., X_n]\) 在 \(k(X_1, ..., X_n)\) 的有限拟伽罗瓦扩张 K 中的整闭包的情形。令 K' 为 K 中 \(k(X_1, ..., X_n)\) 的纯不可分闭包，A' 为 \(k[X_1, ..., X_n]\) 在 K' 中的整闭包。利用习题 17, b)，归约到为 A' 证明所要求的性质。仿照 V, § 3, no. 2 的定理 2 的证明，将 K' 嵌入扩张 \(k'(X_1^{q^{-1}}, ..., X_n^{q^{-1}}) = K''\) 中，其中 \(k'\) 是 \(k\) 的纯不可分扩张，并依据习题 17, a) 归约到 A′ 是 \(k[X_1, ..., X_n]\) 在 K'' 中的整闭包的情形。然后注意到 A' = \(k'[X_1^{q^{-1}}, ..., X_n^{q^{-1}}]\)（V, § 1, No. 3, 命题 13 的推论 2），且 A' 是正则环（p. 94, 习题 6）。

\begin{enumerate}
\def\labelenumi{\arabic{enumi})}
\setcounter{enumi}{18}
\item
  设 A 为域 k 上的有限型整环。满足 \(m \in \text{Specmax}(A)\) 且 \(A_{m}\) 为正则环的极大理想集合在 \(\text{Specmax}(A)\) 中开且稠密。
\item
  设 \(k\) 为一个域，\(k_{i}\)（\(i = 1, 2\)）为 \(k\) 的两个扩张；若超越次数有限，则令 \(n_{i}\)（\(i = 1, 2\)）为 \(k_{i}\) 在 \(k\) 上的超越次数，否则令其为符号 +\(\infty\)。证明有
\end{enumerate}

\[
\dim (k _ {1} \otimes_ {k} k _ {2}) = \inf (n _ {1}, n _ {2})
\]

并且 \(k_{1} \otimes_{k} k_{2}\) 是诺特环，当扩张 \(k_{i}\) 中至少有一个是有限型时。

\begin{enumerate}
\def\labelenumi{\arabic{enumi})}
\setcounter{enumi}{20}
\item
  设 I 为定向有序集，\((\mathbf{A}_{i}, \varphi_{ij})\) 为以 A 为归纳极限的环的归纳系统。假设对于每个 i，环 \(A_{i}\) 都是诺特环；对于每个对 \((i, j)\)（其中 i \textless{} j），\(\varphi_{ij}\) 使 \(A_{j}\) 成为平坦的 \(A_{i}\)-模；且 A 是诺特环。证明若 \(A_{i}\) 对所有 \(i \in I\) 都是正则环，则 A 是正则环。（可以先归约到 \(A_{i}\) 和 A 都是局部环且 \(\varphi_{ij}\) 是局部同态的情形。然后使用 p.94 习题 3 和 p.96 习题 14。）
\item
  \begin{enumerate}
  \def\labelenumii{\alph{enumii})}
  \tightlist
  \item
    设 \(k\) 为一个域，A 为有限型 \(k\)-代数，\(k'\) 为 \(k\) 的有限可分扩张。证明若 A 为正则环，则 \(A \otimes_k k'\) 为正则环。（可以使用 p.96 习题 14, b)。）
  \end{enumerate}
\end{enumerate}

\begin{enumerate}
\def\labelenumi{\alph{enumi})}
\setcounter{enumi}{1}
\tightlist
\item
  在 A 具有相同假设的条件下，设 \(k'\) 为 \(k\) 的有限纯不可分扩张。证明一般而言 \(A \otimes_k k'\) 不是正则环（取 \(A = k'\)）。
\end{enumerate}

\begin{enumerate}
\def\labelenumi{\arabic{enumi})}
\setcounter{enumi}{22}
\tightlist
\item
  设 \(k\) 为特征指数为 \(p\) 的域，A 为有限型 \(k\)-代数，\(\mathfrak{p}\) 为 A 的素理想。证明下列条件等价：
\end{enumerate}

\begin{enumerate}
\def\labelenumi{(\roman{enumi})}
\item
  环 \(A_{\mathfrak{p}} \otimes_{k} k'\) 对 k 的每个扩张 \(k'\) 都是正则环；
\item
  环 \(A_{\mathfrak{p}} \otimes_{k} k^{p^{-\infty}}\) 是正则环；
\item
  环 \(\mathbf{A}_{\mathfrak{p}} \otimes_{k} k'\) 对每个有限且纯不可分的扩张 \(k'\)（扩张 \(k\)）都是正则环。
\end{enumerate}

 （为证明 (ii) \(\Rightarrow\) (iii)，使用 p.96 习题 14, a)。为证明 (iii) \(\Rightarrow\) (i)，只需证明 \(A_{\mathfrak p} \otimes_k k'\) 在 \(k'\) 是 \(k\) 的有限生成扩张时为正则环（习题 21），并且当 \(k'\) 是纯超越扩张的纯不可分扩张时（习题 22 和 14, a)），且被纯超越扩张的一个有限拟伽罗瓦扩张所控制时也是如此。然后把 \(k'\) 控制在扩张 \(k''(\mathrm{X}_1, ..., \mathrm{X}_n)\) 中，其中 \(k''\) 是 \(k\) 的有限纯不可分扩张，并归约到 \(k' = k''(\mathrm{X}_1, ..., \mathrm{X}_n)\) 的情形。由习题 14, a)，置 \(\mathbf{B} = \mathbf{A}_{\mathfrak{p}} \otimes_k k''\)，根据 (iii) 得到一个正则代数，而 \(\mathbf{A}_{\mathfrak{p}} \otimes_k k' = \mathbf{B} \otimes_{k''} k''(\mathbf{X}_1, ..., \mathbf{X}_n)\) 根据 p.94 习题 6, e) 是正则代数。）

若 \(p \in \text{Spec}(A)\) 具有上述等价性质，则称 p 为光滑点。

\begin{enumerate}
\def\labelenumi{\arabic{enumi})}
\setcounter{enumi}{23}
\item
  设 \(k\) 为一个域，A 为 \(k\)-代数，\(k'\) 为 \(k\) 的纯不可分扩张，且 \(\mathbf{A}' = \mathbf{A} \otimes_k k'\)。证明典范映射 \(\operatorname{Spec}(\mathbf{A}') \to \operatorname{Spec}(\mathbf{A})\) 是同胚。（当 A 为域时，\(\mathbf{A}'\) 的维数为 0（习题 20）；证明它只有一个素理想。由此可推出，对于 A 的每个素理想 \(\mathfrak{p}\)，\(\mathfrak{p}'\) 是 \(\mathfrak{p}\mathbf{A}'\) 的根基，且映射 \(\mathfrak{p} \mapsto \mathfrak{p}'\) 是连续的。）
\item
  设 \(k\) 为一个域，A 为有限型 \(k\)-代数，\(\operatorname{Lis}(A)\) 为由光滑的 \(\mathfrak{p} \in \operatorname{Spec}(A)\) 构成的集合（习题 23）。证明 \(\operatorname{Lis}(A)\) 在 \(\operatorname{Spec}(A)\) 中是开集。（使用习题 18、23 和 24。）
\item
  设 \(k\) 为一个域，A 为有限型整环 \(k\)-代数。证明 Lis(A) 稠密当且仅当 A 是可分的 \(k\)-代数；等价地，A 的分式域是 \(k\) 的可分扩张（A, V, p.115）。
\item
  设 \(k\) 为一个域，A 为有限型且整的可分 \(k\)-代数。A 的光滑极大理想集合是 Specmax(A) 的开稠密子集。
\item
  设 \(k\) 为一个域，\(k'\) 为 \(k\) 的扩张。证明对于每个扩张 \(k''\)（它是 \(k\) 的扩张），\(k' \otimes_k k''\) 是正则诺特环当且仅当 \(k'\) 是 \(k\) 的有限型可分扩张。
\end{enumerate}

29) 设 k 为一个域，A 为完备诺特局部 k-代数，\(\kappa_{A}\) 为 A 的剩余域。

\begin{enumerate}
\def\labelenumi{\alph{enumi})}
\item
  假设 \(\kappa_{\mathbf{A}}\) 是 \(k\) 的可分扩张。证明存在 \(k\)-代数同态，从 \(\kappa_{\mathbf{A}}\) 到 A，它是从 A 到 \(\kappa_{\mathbf{A}}\) 的典范同态的截面（IX, § 3, n° 2, 命题 1）。
\item
  假设 \(\kappa_{\mathbf{A}}\) 是 \(k\) 的纯不可分扩张。证明不一定存在 \(k\)-代数截面，从 \(\kappa_{\mathbf{A}}\) 到 A。（取特征为 2 的域 \(k\)，取不是平方的元素 \(a\)，它属于 \(k\)（例如 \(k = \mathbf{F}_2(\mathrm{T})\)，\(a = \mathrm{T}\)），取理想 \(\mathfrak{p}\)，它是 \(k[\mathrm{X}]\) 中由 \(\mathrm{X}^2 + a\) 生成的理想，并置 \(\mathrm{A} = \widehat{k[\mathrm{X}]_{\mathfrak{p}}}\)。）
\item
  假设 A 为正则环，且 \(\kappa_{\mathbf{A}}\) 是 \(k\) 的可分扩张。证明 A 作为 \(k\)-代数同构于 \(\kappa_{\mathbf{A}}[[\mathrm{T}_{1}, ..., \mathrm{T}_{n}]]\)，其中 \(n = \dim(\mathbf{A})\)。
\item
  证明当 \(\kappa_{\mathrm{A}}\) 不是 \(k\) 的可分扩张时，情形不一定如此。
\end{enumerate}

\begin{enumerate}
\def\labelenumi{\arabic{enumi})}
\setcounter{enumi}{29}
\item
  设 \(k\) 为一个域，A 为局部、诺特且完备的 \(k\)-代数。若对于每个有限扩张 \(k'\)（它是 \(k\) 的扩张），环 \(A \otimes_k k'\) 都是正则环，则称 A 形式光滑。证明 A 形式光滑当且仅当对于每个有限纯不可分扩张 \(k'\)（它是 \(k\) 的扩张），环 \(A \otimes_k k'\) 都是正则环（可以使用 p.96 习题 14，并仿照 p.97 习题 23 的证明）。
\item
  设 \(k\) 为一个域。
\end{enumerate}

\begin{enumerate}
\def\labelenumi{\alph{enumi})}
\item
  设 A 为诺特、局部、完备且正则的 \(k\)-代数，其剩余域 \(\kappa_{\mathrm{A}}\) 是 \(k\) 的可分扩张。证明 A 形式光滑且同构于 \(\kappa_{\mathrm{A}}[[\mathbf{T}_{1}, ..., \mathbf{T}_{n}]]\)。
\item
  设 B 为有限型 \(k\)-代数，\(\mathfrak{p} \in \operatorname{Spec}(\mathbf{B})\) 为光滑点（p.97, 习题 23）。证明 \(\hat{\mathbf{B}}_{\mathfrak{p}}\) 是形式光滑的 \(k\)-代数。
\item
  设 \(\kappa\) 为 \(k\) 的有限生成扩张，且 \(n\in\mathbf{N}\)。证明存在形式光滑的 \(k\)-代数，其剩余域为 \(\kappa\)，维数为 \(n\)。
\item
  设 A 为形式光滑的 k-代数，其剩余域 \(\kappa_{A}\) 不是 k 的可分扩张。证明 A 作为 k-代数不同构于 \(\kappa_{A}[[T_{1},...,T_{n}]]\)。
\end{enumerate}

32) 设 \(k\) 为一个域，K 为 \(k\) 的有限生成扩张，\(n\) 为正整数，A、B 为两个形式光滑 \(k\)-代数，维数为 \(n\)，且 \(\kappa_{A}\) 和 \(\kappa_{B}\) 作为 \(k\) 的扩张都同构于 K。拟证明 A、B 是同构的 \(k\)-代数。

\begin{enumerate}
\def\labelenumi{\alph{enumi})}
\item
  归约到 K 是 k 的有限纯不可分扩张的情形（习题 29, a)）。
\item
  考虑 \(C = A \hat{\otimes}_{k} B\)，它是 \(A \otimes_{k} B\) 关于理想的完备化代数
\end{enumerate}

\[
\left(\mathbf {A} \otimes_ {k} \mathfrak {m} _ {\mathbf {B}}\right) + \left(\mathfrak {m} _ {\mathbf {A}} \otimes_ {k} \mathbf {B}\right).
\]

证明 C 是诺特、局部、完备且正则的 \(k\)-代数。（为证明最后一点，使用从 A 到 C 的典范同态 \(\rho\)。我们将证明它使 C 成为平坦的 A-模，且 \(C \otimes_{\mathbf{A}} \kappa_{\mathbf{A}}\) 同构于 \(\kappa_{\mathbf{A}} \otimes_{k} \mathbf{B}\)。然后使用 p.96 习题 14。）

\begin{enumerate}
\def\labelenumi{\alph{enumi})}
\setcounter{enumi}{2}
\tightlist
\item
  证明 \(\rho\) 在剩余域上诱导同构，并将 \(m_{\mathbf{A}}/m_{\mathbf{A}}^{2}\) 单射到 \(m_{\mathbf{C}}/m_{\mathbf{C}}^{2}\) 中。然后构造 \(\pi\) 作为 \(\rho\) 的一个回缩，它是 \(k\)-代数同态。证明 \(\pi \circ \rho': B \to A\) 是 \(k\)-代数同构。（我们将从 B 到 C 的典范同态记为 \(\rho'\)。）
\end{enumerate}

